\documentclass[11pt]{article}
\usepackage[a4paper,margin=2.3cm]{geometry}
\usepackage{amsmath,amssymb}
\usepackage{graphicx}
\usepackage{booktabs}
\usepackage{xcolor}
\usepackage{titlesec}
\usepackage{caption}
\usepackage{hyperref}
\usepackage{array}
\usepackage{enumitem}
\usepackage{fancyhdr}

\definecolor{edhecred}{RGB}{158,27,50}
\definecolor{edhecgray}{RGB}{77,77,77}

\titleformat{\section}{\large\bfseries\color{edhecred}}{\thesection}{0.6em}{}
\titleformat{\subsection}{\normalsize\bfseries\color{edhecgray}}{\thesubsection}{0.6em}{}
\hypersetup{colorlinks=true, linkcolor=edhecred, urlcolor=edhecred, citecolor=edhecred}

\graphicspath{{outputs/}{python_labs/outputs/}}

\pagestyle{fancy}
\fancyhf{}
\lhead{\small Advanced Fixed Income and Credit --- Group Assignment}
\rhead{\small\thepage}
\renewcommand{\headrulewidth}{0.3pt}

\newcommand{\question}[1]{\subsection*{\textcolor{edhecred}{#1}}}

\begin{document}

%====================================================================
\begin{titlepage}
\centering
\vspace*{2cm}
{\LARGE\bfseries\color{edhecred} Advanced Fixed Income and Credit}\\[0.4cm]
{\Large Group Assignment}\\[1.2cm]
{\large Part I: Deutsche Bank --- Finding Relative-Value Trades (HBS 205-059)}\\
{\large Part II: Python Labs --- US Treasury Yield Curve Reaction to Trump Tariff Threats}\\[2cm]
{\large EDHEC Business School --- MSc\\ Advanced Fixed Income and Credit, Thibault Lair}\\[2cm]
{\large \today}\\[1cm]
\vfill
{\small \textit{Group members: [FILL IN NAMES HERE]}}
\end{titlepage}

\tableofcontents
\newpage

%====================================================================
\section{Part I --- The Deutsche Bank Case Study}

\subsection*{Setup and conventions}
Settlement date: \textbf{August 15, 2003}. Exhibit 1 lists 52 Treasury coupon
bonds whose maturities fall on successive semi-annual dates (Feb 15 / Aug 15)
from 2004 to 2029 with no gaps, i.e.\ exactly one bond at each of the nodes
$t=0.5,1.0,\dots,26.0$ years. This allows the classic \emph{successive
substitution} bootstrap: because each bond introduces exactly one new unknown
discount factor, we solve sequentially from the shortest to the longest
maturity.

\textbf{Compounding convention.} Exhibit 4's footnote states explicitly that
Deutsche Bank's own model yields are \emph{Bond-Equivalent Yields} (BEY),
i.e.\ twice the semi-annual compounded yield. Since U.S.\ Treasuries pay
semi-annual coupons, we derive our own zero-coupon yields in the identical
BEY convention throughout, so that every comparison against the DB model in
Q1.2 is apples-to-apples (a continuously-compounded convention would have
introduced a small but systematic bias in one direction).
$$\text{BEY} = 2\left(D^{-1/k}-1\right), \qquad D = \text{discount factor to period } k \text{ (6-month steps)}.$$

%--------------------------------------------------------------------
\question{Q1.1 --- Bootstrapping the discount curve}

\textbf{Methodology.} For bond $k$ (maturing at period $k$, coupon $c$ per
annum, price $P_k$):
$$P_k = \frac{c}{2}\sum_{j=1}^{k-1} D_j + \left(\frac{c}{2}+100\right)D_k
\quad\Longrightarrow\quad
D_k = \frac{P_k - \frac{c}{2}\sum_{j=1}^{k-1}D_j}{\frac{c}{2}+100}.$$
Solved in order of increasing maturity, using only already-known $D_j$, $j<k$.

\textbf{Validation.} Our bootstrapped curve matches Deutsche Bank's own
Exhibit 4 model yields \emph{exactly} at 2Y (1.873\%) and 10Y (4.711\%) ---
precisely the maturities the case identifies as DB's own calibration points
for its three-factor model. This is strong internal validation that both the
bootstrap algorithm and the BEY convention choice are correct.

\begin{table}[h]
\centering
\small
\begin{tabular}{rrrr}
\toprule
Maturity (y) & Discount Factor & Zero Yield (BEY) \\
\midrule
0.5  & 0.99561 & 0.882\% \\
1.0  & 0.98818 & 1.193\% \\
2.0  & 0.96341 & 1.873\% \\
5.0  & 0.84156 & 3.480\% \\
10.0 & 0.62775 & 4.711\% \\
15.0 & 0.42935 & 5.717\% \\
20.0 & 0.30947 & 5.951\% \\
25.0 & 0.23299 & 5.913\% \\
26.0 & 0.22024 & 5.905\% \\
\bottomrule
\end{tabular}
\caption{Bootstrapped zero-coupon curve, selected maturities (full 52-point curve in \texttt{q1\_bootstrapped\_curve.csv}).}
\end{table}

The curve shape is the standard concave, upward-sloping Treasury curve of
the period, flattening from about 15Y onward and showing a very slight
downward drift between 20Y and 26Y.

%--------------------------------------------------------------------
\question{Q1.2 --- Comparison with Deutsche Bank's model, trading signals}

\textbf{Logic.} Yield and price move in opposite directions. If our
bootstrapped (market-implied) yield is \emph{higher} than DB's model yield,
the bond's market price is below the model's ``fair'' price $\Rightarrow$ the
bond is \textbf{cheap} $\Rightarrow$ \textbf{BUY}. If the market yield is
\emph{lower}, the bond is \textbf{rich} $\Rightarrow$ \textbf{SELL}.

\begin{table}[h]
\centering
\small
\begin{tabular}{rrrrl}
\toprule
Maturity (y) & DB Model Yield & Market Yield (Q1.1) & Diff (bps) & Signal \\
\midrule
1  & 1.2443\% & 1.1930\% & $-5.13$ & SELL (rich) \\
2  & 1.8727\% & 1.8727\% & $+0.00$ & fair (DB calibration point) \\
3  & 2.4110\% & 2.4711\% & $+6.01$ & \textbf{BUY (cheap)} \\
4  & 2.9665\% & 3.0084\% & $+4.19$ & \textbf{BUY (cheap)} \\
5  & 3.4454\% & 3.4799\% & $+3.45$ & \textbf{BUY (cheap)} \\
6  & 3.8557\% & 3.8817\% & $+2.60$ & \textbf{BUY (cheap)} \\
7  & 4.1996\% & 4.2203\% & $+2.07$ & \textbf{BUY (cheap)} \\
8  & 4.4677\% & 4.4865\% & $+1.88$ & \textbf{BUY (cheap)} \\
9  & 4.6528\% & 4.6640\% & $+1.12$ & \textbf{BUY (cheap)} \\
10 & 4.7107\% & 4.7107\% & $-0.00$ & fair (DB calibration point) \\
15 & 5.7160\% & 5.7167\% & $+0.07$ & fair \\
20 & 5.9517\% & 5.9513\% & $-0.04$ & fair \\
25 & 5.9315\% & 5.9128\% & $-1.87$ & SELL (rich) \\
\bottomrule
\end{tabular}
\caption{Market vs. DB model comparison and trading signals.}
\end{table}

\textbf{Pattern.} The belly of the curve (3--9 years) is uniformly
\emph{cheap} relative to the DB model, by 1--6 bps, suggesting a BUY
candidate basket in that region. The very short end (1Y) and the long end
(25Y) look \emph{rich}. The 2Y, 10Y, 15Y and 20Y points are essentially fair
by construction or coincidence (2Y and 10Y are DB's own calibration anchors).

%--------------------------------------------------------------------
\question{Q1.3 --- Is this risk-free arbitrage?}

No. Although the strategy is directional and model-based, it is \textbf{not}
a riskless arbitrage, for several reasons:

\begin{itemize}[leftmargin=1.2em]
\item \textbf{Model risk.} The ``mispricing'' is only a mispricing relative
to Deutsche Bank's own three-factor curve model. If the model is
mis-specified (wrong number of factors, wrong functional form, stale
calibration), the signal is simply wrong --- there is no model-independent
proof that the belly of the curve is ``truly'' cheap.
\item \textbf{Convergence/timing risk.} Even if the model is correctly
specified on average, there is no guarantee \emph{when}, or even
\emph{whether}, market yields converge back toward the model yields within
the trader's holding horizon. The position carries interim mark-to-market
risk.
\item \textbf{No riskless, capital-structure-free payoff.} A true arbitrage
requires a replicating, self-financing portfolio with a knowable, riskless
terminal payoff (e.g.\ a textbook cash-and-carry). Here the ``trade'' is an
outright directional bet on convergence, funded at some financing rate, with
basis risk between the position and any hedge.
\item \textbf{Curve risk while waiting.} Between trade inception and
convergence, the \emph{whole} curve can move (parallel shifts, twists), and
unless the position is duration/key-rate hedged, those moves swamp the small
model-implied edge of a few basis points.
\end{itemize}
In short, this is a \textbf{relative-value} trade: a model-informed bet with
positive expected value under the trader's own beliefs about the model's
validity, not a riskless arbitrage in the classical sense.

%--------------------------------------------------------------------
\question{Q2.1 --- Cubic spline of the zero-coupon curve}

\textbf{Methodology.} Natural cubic spline (zero second derivative at the
boundary knots) through the five knots specified (2, 5, 10, 15, 20 years),
fit directly on the Q1.1 bootstrapped BEY yields. Outside $[2,20]$y the
spline is extrapolated using its boundary slope.

\begin{figure}[h]
\centering
\includegraphics[width=0.85\textwidth]{q2_spline_vs_bootstrap.png}
\caption{Cubic spline (red) vs. raw bootstrapped zero yields (grey dots). Shaded regions are extrapolated.}
\end{figure}

\textbf{Fit quality.} Within the knot range $[2,20]$y: mean residual
$=-5.50$\,bps, std $=8.28$\,bps, max $|$residual$| = 29.91$\,bps. The
residuals are \emph{not} randomly scattered: the largest errors cluster
sharply between 7.5 and 13 years (up to $-30$\,bps at 11.5Y), exactly the
maturity window in which Exhibit 1 contains several legacy, high-coupon
bonds issued in the early 1980s (coupons of 13.25\%, 12.5\%, 11.25\%,
10.625\%). These bonds plausibly carry idiosyncratic pricing noise (age,
liquidity, possible embedded optionality/callability common to that
vintage), which the raw bootstrap faithfully reproduces pointwise but which
the smooth spline, by construction, averages away.

%--------------------------------------------------------------------
\question{Q2.2 --- Does using the spline change the Q1.2 conclusions?}

We repeat the Q1.2 comparison using the \textbf{spline-implied} yield
(instead of the raw bootstrap) at every maturity used in the original
table, since the spline value differs from the raw bootstrap except
exactly at the five knots (2,5,10,15,20).

\begin{table}[h]
\centering
\small
\begin{tabular}{rrrrll}
\toprule
Maturity (y) & DB Model & Spline Yield & Diff (bps) & Spline Signal & \textit{Changed vs. Q1.2?} \\
\midrule
1  & 1.2443\% & 1.2846\% & $+4.03$  & BUY  & \textbf{YES} (was SELL) \\
2  & 1.8727\% & 1.8727\% & $+0.02$  & fair & no (knot) \\
3  & 2.4110\% & 2.4608\% & $+4.98$  & BUY  & no \\
4  & 2.9665\% & 3.0096\% & $+4.31$  & BUY  & no \\
5  & 3.4454\% & 3.4799\% & $+3.45$  & BUY  & no (knot) \\
6  & 3.8557\% & 3.8439\% & $-1.18$  & SELL & \textbf{YES} (was BUY) \\
7  & 4.1996\% & 4.1195\% & $-8.01$  & SELL & \textbf{YES} (was BUY) \\
8  & 4.4677\% & 4.3363\% & $-13.14$ & SELL & \textbf{YES} (was BUY) \\
9  & 4.6528\% & 4.5235\% & $-12.93$ & SELL & \textbf{YES} (was BUY) \\
10 & 4.7107\% & 4.7107\% & $-0.00$  & fair & no (knot) \\
15 & 5.7160\% & 5.7167\% & $+0.07$  & fair & no (knot) \\
20 & 5.9517\% & 5.9513\% & $-0.04$  & fair & no (knot) \\
25 & 5.9315\% & 6.1859\% & $+25.44$ & BUY  & \textbf{YES} (was SELL) \\
\bottomrule
\end{tabular}
\caption{Q1.2 signals recomputed off the Q2.1 spline curve.}
\end{table}

\textbf{This matters a great deal.} Six of the thirteen signals (1, 6, 7, 8,
9, 25 years) \emph{flip sign} once the spline replaces the raw bootstrap.
Two distinct mechanisms are at work:

\begin{enumerate}[leftmargin=1.4em]
\item \textbf{Extrapolation risk (1Y and 25Y).} These two points lie
\emph{outside} the spline's knot range $[2,20]$y. The natural spline's
boundary-slope extrapolation is a pure modelling artifact with no claim to
describing the true curve shape beyond the data used to fit it, so any
signal built on it is unreliable almost by definition.
\item \textbf{Belly-of-curve smoothing (6--9Y).} This is the more
substantive finding. The original Q1.2 ``cheap'' signal in 6--9Y was driven,
at least in part, by the \emph{same} idiosyncratic legacy-bond noise
identified in Q2.1's residual analysis. Once the spline smooths through
that local bump, the belly of the curve actually looks \emph{rich}, not
cheap, relative to the DB model.
\end{enumerate}

\textbf{Conclusion.} The trading conclusion from Q1.2 is \emph{not robust}
to the choice of curve-fitting methodology in exactly the region (6--9Y)
where the raw bootstrap is most likely contaminated by non-fundamental,
bond-specific noise. A prudent trader would treat the raw-bootstrap BUY
signal in 6--9Y with real skepticism, and would want to understand
\emph{why} those bonds are mispriced (liquidity? embedded option? stale
quote?) before trading on it, rather than assuming genuine relative value.

%--------------------------------------------------------------------
\question{Q3.1 --- Nelson-Siegel-Svensson fit}

\textbf{Methodology.} Jointly estimating $(\beta_0,\beta_1,\beta_2,\beta_3,
\tau_1,\tau_2)$ by unconstrained nonlinear least squares proved numerically
fragile: this curve is concave with a single mild hump and no second hump,
so $\beta_0$ (level) and $\beta_2$ (curvature) become highly collinear once
$\tau_1$ is pushed above $\sim$10 years (the curvature loading then barely
decays within our 26-year range). Different, equally-well-fitting local
optima emerged from different starting points.

We instead use the standard \textbf{Diebold--Li two-step procedure}: for
each candidate $(\tau_1,\tau_2)$ on a fine grid, the loadings become fixed
basis functions of maturity, so the betas enter \emph{linearly} and are
solved by OLS --- a well-posed, unique, convex problem for every candidate
$\tau$. We then pick the $(\tau_1,\tau_2)$ minimizing total RSS.

\begin{figure}[h]
\centering
\includegraphics[width=0.78\textwidth]{q3_nss_fit.png}
\caption{Nelson-Siegel fit vs. the Q1.1 bootstrapped curve.}
\end{figure}

\textbf{Model selection.} The plain 3-factor Nelson-Siegel achieves RMSE
$=7.96$ bps ($\beta_0=0.695\%$, $\beta_1=-0.140\%$, $\beta_2=17.926\%$,
$\tau_1=12.08$y). The full 4-factor Svensson only improves RMSE to $7.28$
bps (an improvement of $0.68$ bps), below our pre-set materiality threshold
of 1.5 bps. We therefore adopt the \textbf{simpler, better-identified
3-factor Nelson-Siegel model}: the data show no evidence of a genuine
second hump, so the extra parameter and its attendant identification
problems are not worth the marginal fit improvement.

\begin{table}[h]
\centering
\small
\begin{tabular}{lr}
\toprule
Parameter & Value \\
\midrule
$\beta_0$ (level)     & 0.695\% \\
$\beta_1$ (slope)     & $-0.140\%$ \\
$\beta_2$ (curvature) & 17.926\% \\
$\tau_1$              & 12.081 years \\
RMSE                  & 7.96 bps \\
\bottomrule
\end{tabular}
\caption{Final adopted Nelson-Siegel parameters.}
\end{table}

%--------------------------------------------------------------------
\question{Q3.2 --- Factor loadings and economic drivers}

\begin{figure}[h]
\centering
\includegraphics[width=0.78\textwidth]{q3_nss_loadings.png}\\[0.3cm]
\includegraphics[width=0.78\textwidth]{q3_nss_contributions.png}
\caption{Top: factor loadings. Bottom: factor contributions ($\beta_i \times$ loading$_i$).}
\end{figure}

The contribution decomposition reveals a striking feature: the
\textbf{curvature} factor ($\beta_2=17.9\%$) does almost all of the
curve-shaping work, rising smoothly from 0 at $m=0$ to its asymptote around
20 years, while the level contribution is small and essentially flat
($0.695\%$) and the slope contribution is tiny and nearly flat
($-0.140\%$). This is counter-intuitive at first glance --- in textbook
Nelson-Siegel intuition, ``level'' should dominate a curve of this overall
magnitude --- but it is a direct, mechanical consequence of $\tau_1$ being
large ($12.08$ years) relative to our 26-year maturity range: the curvature
loading function $\left[\frac{1-e^{-m/\tau_1}}{m/\tau_1}-e^{-m/\tau_1}\right]$
simply has not decayed back to zero by $m=26$, so it is statistically
\emph{behaving like} a second, slowly-varying level/slope factor over the
range of our data, and the OLS step naturally loads most of the explanatory
power onto it rather than onto $\beta_0$ precisely because $\beta_0$'s
loading is flat and uninformative about curve \emph{shape}, whereas the
curvature loading can trace the whole concave shape on its own.

Economically, this echoes Deutsche Bank's own Exhibit 3 three-factor model
description: a persistent, long-horizon ``long-rate''/level-type factor
that evolves slowly is naturally captured by a Nelson-Siegel factor whose
decay parameter is large relative to the fitting window, rather than by the
model's nominal level factor, which by construction carries no shape
information at all. In other words, when the sample's maturity range is not
much longer than $\tau_1$, Nelson-Siegel's ``curvature'' factor and a
slow-moving macro ``long rate'' factor become difficult to tell apart ---
exactly the level/curvature identifiability issue we had already flagged
during estimation (Q3.1).

%--------------------------------------------------------------------
\question{Q4.1 --- Shocking the 10Y key rate}

\textbf{Methodology.} We work off the Q2.1 \emph{spline} curve (not the raw
bootstrap or NSS), because a key-rate shock is only a well-defined operation
relative to a smooth curve model with discrete nodes --- and the spline's
knots (2,5,10,15,20y) are exactly the assignment's key rates. Each key-rate
shock is the standard Ho (1992) triangular (``tent'') perturbation: full
size at its own knot, decaying linearly to zero at the two adjacent knots,
and held flat beyond the outer knots. We verified by construction that the
five key-rate shocks sum to exactly a parallel shift at every maturity.

\begin{figure}[h]
\centering
\includegraphics[width=0.85\textwidth]{q4_1_shock_10y.png}
\caption{Left: base vs. 10Y-shocked curve. Right: the shock itself, a clean triangular ``tent'', zero outside $[5,20]$y.}
\end{figure}

A $+50$bps shock to the 10Y key rate raises yields in a triangular tent
shape centered at 10Y: zero impact at the adjacent 5Y and 20Y knots, ramping
linearly to the full $+50$bps exactly at 10Y. No maturity below 5Y or above
20Y is affected at all --- this strict locality is the defining feature that
distinguishes a key-rate shock from a parallel shift.

%--------------------------------------------------------------------
\question{Q4.2 --- Portfolio KRD to the 10Y shock}

\textbf{Portfolio construction.} \$1,000,000 split equally by \emph{dollar
market value} across all 52 bonds ($\approx\$19{,}230.77$ of market value
in each), translated into a face-value holding of
$(\text{dollar amount}/\text{market price})\times 100$ per bond.

Repricing consistently off the spline curve (base case and shocked case
both use the spline-implied yields, since the whole point of a key-rate
exercise is to measure the portfolio's sensitivity to the curve
\emph{model}):

\begin{itemize}
\item Base ``model NAV'' (spline-repriced): \$1,002,389.09 (vs. the
\$1,000,000 target --- the \$2,389.09 gap is exactly the Q2.2 spline
fitting residual, not a hedging error).
\item Shocked NAV (10Y $+50$bps): \$992,156.79.
\end{itemize}
$$\boxed{\text{KRD}_{10Y} = \$992{,}156.79 - \$1{,}002{,}389.09 = -\$10{,}232.30}$$
i.e.\ $-102.08$ bps of NAV per 50bps shock at the 10Y key rate. The negative
sign is exactly as expected: a long bond portfolio loses value when yields
rise.

%--------------------------------------------------------------------
\question{Q4.3 --- Full key-rate duration profile}

\begin{table}[h]
\centering
\small
\begin{tabular}{rrr}
\toprule
Key Rate & KRD (\$) & KRD (\% of NAV) \\
\midrule
2Y  & $-2{,}640.21$  & $-0.263\%$ \\
5Y  & $-7{,}184.23$  & $-0.717\%$ \\
10Y & $-10{,}232.30$ & $-1.021\%$ \\
15Y & $-8{,}305.04$  & $-0.829\%$ \\
20Y & $-10{,}537.50$ & $-1.051\%$ \\
\bottomrule
\end{tabular}
\caption{Key-rate duration profile of the \$1M equally-weighted portfolio.}
\end{table}

\begin{figure}[h]
\centering
\includegraphics[width=0.7\textwidth]{q4_3_krd_profile.png}
\caption{KRD profile: the portfolio's dollar P\&L sensitivity to each key rate.}
\end{figure}

\textbf{Internal consistency check.} The sum of the five individual KRDs is
$-\$38{,}899.28$. The KRD of a single $+50$bps \emph{parallel} shock (since
the five triangular shocks sum exactly to a parallel shift) is
$-\$38{,}619.06$. These match to within \$280.22 ($\sim$0.03\% of NAV), the
small residual coming from the convexity of bond pricing in yield --- KRDs
are summed \emph{linearly} but repricing is nonlinear in yield, so a
sum-of-parts never \emph{exactly} equals the whole for a non-trivial shock
size.

\textbf{Interpretation.} All five KRDs are negative, as expected for a long
portfolio. The profile is heaviest at 10Y and 20Y ($\approx -\$10.2$k and
$-\$10.5$k), and lightest at the short end (2Y, $-\$2.6$k) --- consistent
with the portfolio holding roughly equal dollar amounts at every maturity,
so the longer-duration bonds (with more distant, more convexity-sensitive
cash flows) dominate the risk profile even though every bond carries the
same market-value weight.

\newpage
%====================================================================
\section{Part II --- Python Labs: US Treasury Yield Curve Reaction to Trump Tariff Threats}

\subsection*{Data}
US Treasury yields (3M to 30Y) fetched from FRED (DGS series), 2015-01-02 to
2025-08-01 (2{,}647 trading days). Event dates: the assignment's own
\texttt{trump\_actions} dictionary (tariff threats/actions, both
administrations) and \texttt{country\_responses} dictionary (foreign
retaliation announcements), merged into a single 58-event table (we did not
have access to non-US yield data, so ``reaction of targeted countries'' is
operationalized as the \emph{US} Treasury curve's reaction specifically on
retaliation-announcement dates, as a complement to the reaction on Trump's
own action dates --- this assumption is stated explicitly as the data
available to us is US-only).

\question{Q1.1 --- Event-study methodology}
For each event date $t_0$ and each horizon $d\in[-30,+30]$ trading days, we
compute the \textbf{raw change} $y_{t_0+d}-y_{t_0}$, averaged across all 58
events. We then net out ``normal'' yield drift unrelated to the events by
computing the identical same-horizon change over all \textbf{non-event}
dates (any date more than 30 trading days from every event), and define
$$\text{Abnormal change}(d) = \overline{\Delta y}^{\text{event}}(d) - \overline{\Delta y}^{\text{non-event}}(d).$$
This is applied to the raw 10Y yield and to three constructed composite
series: \textbf{level} (10Y), \textbf{slope} (10Y$-$2Y), and
\textbf{curvature} ($2\times$5Y $-$ 2Y $-$ 10Y).

\question{Q1.2 --- Results: level, slope and curvature}

\begin{figure}[h]
\centering
\includegraphics[width=0.85\textwidth]{es_10y_abnormal.png}
\caption{10Y yield: event-conditional mean change, non-event baseline, and abnormal change.}
\end{figure}

\begin{table}[h]
\centering
\small
\begin{tabular}{rrrr}
\toprule
Horizon (trading days) & 10Y / Level (bps) & Slope (bps) & Curvature (bps) \\
\midrule
$-30$ & $+11.46$ & $+1.30$ & $+2.90$ \\
$-10$ & $+7.26$  & $+2.36$ & $+1.50$ \\
$-5$  & $+3.40$  & $+1.22$ & $+0.41$ \\
$-1$  & $+0.56$  & $-0.49$ & $+0.23$ \\
$0$   & $0.00$   & $0.00$  & $0.00$ \\
$+1$  & $+0.38$  & $+0.86$ & $+0.17$ \\
$+5$  & $-0.20$  & $+0.89$ & $-0.33$ \\
$+10$ & $-3.50$  & $+0.24$ & $-1.75$ \\
$+30$ & $-7.72$  & $+2.02$ & $-3.98$ \\
\bottomrule
\end{tabular}
\caption{Abnormal change at selected horizons, all 58 events combined.}
\end{table}

\begin{figure}[h]
\centering
\includegraphics[width=0.95\textwidth]{es_lsc_decomposition.png}
\caption{Abnormal change decomposition: level, slope, curvature.}
\end{figure}

\textbf{Level/curvature} show a large, smooth ``slope'' shape: abnormal
change is strongly positive ($\sim+11$bps) 30 days \emph{before} the event
window and turns negative ($\sim-8$bps) 30 days \emph{after}. \textbf{Slope}
is much smaller in magnitude and does not show this clean monotonic
pattern.

\textbf{Important caveat (anticipating Q2.1).} This smooth, large,
symmetric swing is almost certainly \emph{not} a genuine 60-day reaction to
individual tariff announcements --- a true event effect would show a sharp
jump concentrated right around $d=0$, not a slow 30-day ramp on both sides.
It far more plausibly reflects that tariff events \textbf{cluster in time
periods} (2018--19, 2025) that happened to coincide with broader
macro/rate-cycle regimes (Fed hiking/cutting cycles, QT/QE) in which
Treasury yields were already trending in a particular direction for reasons
having nothing to do with tariffs. This is a textbook event-study
confounding problem, discussed further in Q2.1.

\question{Q1.3 --- Reaction specifically to retaliation (``targeted countries'') events}
Because only US yield data was available to us, we cannot directly observe
how \emph{foreign} yield curves reacted. As a second-best, we can re-run the
identical event-study machinery restricting the event set to only the
\texttt{country\_response} dates (excluding \texttt{trump\_action} dates),
to see whether the \emph{US} curve reacts differently to retaliation news
than to the original US tariff threats. \textit{[This subset run uses the
same \texttt{event\_study()} function filtered to \texttt{event\_type ==
'country\_response'} in place of the full event list --- see the Appendix
code snippet. Once run, insert the resulting abnormal-change table here; the
full-sample result above already nets together both event types.]}

\question{Q1.4 --- First term (2017--2021) vs. second term (2025) comparison}
Similarly, splitting the 58-event sample by the \texttt{term} column
(\texttt{first\_term} vs. \texttt{second\_term}, already tagged in
\texttt{build\_events\_table()}) and re-running the event study on each
subset separately would let us test whether the market's reaction to tariff
threats was systematically different in the two administrations (e.g.\ a
more desensitized/``priced-in'' reaction by 2025 given the extensive 2018--19
precedent). \textit{[Subset run pending --- same function, filtered on
\texttt{events.term=='first\_term'} / \texttt{'second\_term'}; see Appendix.]}

\question{Q2.1 --- Limitations of this event-study design}
\begin{itemize}[leftmargin=1.2em]
\item \textbf{Confounding with the macro regime.} As flagged in Q1.2, events
cluster in specific calendar periods that coincide with broader yield
trends (Fed cycle, QT/QE) unrelated to tariffs; the large, smooth abnormal-
change pattern we find is at least partly a regime effect, not a pure event
effect.
\item \textbf{Overlapping event windows.} With 58 events and 30-day windows
each side, many windows overlap heavily (events often come in bursts of
days or weeks). This violates the independence assumption behind simple
averaging across events and effectively reduces the number of
``independent'' observations well below 58 (pseudo-replication).
\item \textbf{No control for anticipation / leaks.} Several tariff
``threats'' were telegraphed or rumored before their official date, so
markets may have partly priced them in \emph{before} $t_0$, muting the
measured reaction at $t_0$ itself and biasing the pre-event window.
\item \textbf{Confounding with other scheduled macro events.} FOMC meetings,
CPI/NFP releases, and other scheduled announcements fall inside many of
these 60-day windows purely by chance, and their effects are not separated
out.
\item \textbf{Small-sample / multiple-horizon testing.} With 61 horizons
tested per series and no multiple-testing correction, some ``significant-
looking'' abnormal changes at specific horizons could simply be noise.
\item \textbf{Trading-day approximation.} Event dates are snapped to the
next available trading day, and ``business days'' are defined purely by
presence in the FRED series rather than an exact trading calendar,
introducing minor date-alignment imprecision.
\end{itemize}

\appendix
\section*{Appendix: subset-analysis code for Q1.3 / Q1.4}
\begin{verbatim}
country_only = events[events.event_type == 'country_response']['date']
first_term   = events[events.term == 'first_term']['date']
second_term  = events[events.term == 'second_term']['date']

es_country = event_study(ycd['date'], ycd['10Y'], country_only)
es_first   = event_study(ycd['date'], ycd['10Y'], first_term)
es_second  = event_study(ycd['date'], ycd['10Y'], second_term)
\end{verbatim}

\end{document}
